\section{The exact fragment and its representation}
\label{sec:representation}

\subsection{Rational exponential polynomials}

\begin{definition}
Let $\ep_{\Q}$ be the algebra of functions of the form
\begin{equation}
 f(x)=\sum_{\lambda\in\Lambda}P_\lambda(x)e^{\lambda x},
 \qquad \Lambda\subset\Q\text{ finite},\quad P_\lambda\in\Q[x].
 \label{eq:fragment}
\end{equation}
The stored representation combines equal rates, combines equal powers, removes
zero coefficients, and removes zero polynomial blocks.
\end{definition}

A dictionary entry $(\lambda,k)\mapsto c$ denotes $cx^ke^{\lambda x}$.
All three components are exact: $\lambda,c$ are reduced rationals and $k$ is a
nonnegative integer. The zero map denotes the zero function. Polynomial terms
are simply the block with rate zero. Negative rates are first-class; rational
rates such as $1/2$ require no special case.

Addition merges dictionaries. Multiplication uses
\[
 (c x^k e^{\lambda x})(d x^j e^{\mu x})
   =cd\,x^{k+j}e^{(\lambda+\mu)x}.
\]
Reflection, needed for the other half-line, replaces $(\lambda,k,c)$ by
$(-\lambda,k,(-1)^kc)$. Differentiation and a shifted derivative are
\begin{align}
 \D(P_\lambda e^{\lambda x})
   &=(P_\lambda'+\lambda P_\lambda)e^{\lambda x},\\
 (\D-\rho)(P_\lambda e^{\lambda x})
   &=(P_\lambda'+(\lambda-\rho)P_\lambda)e^{\lambda x}.
 \label{eq:shifted-derivative}
\end{align}
No transcendental arithmetic is required to perform any of these operations.

\subsection{Finite-dimensional differential structure}

For a nonzero canonical input define
\[
 d_\lambda=\deg P_\lambda,\qquad
 N=\sum_{\lambda\in\Lambda}(d_\lambda+1),\qquad
 A_f(\D)=\prod_{\lambda\in\Lambda}(\D-\lambda)^{d_\lambda+1}.
\]
The function belongs to the $N$-dimensional space spanned by
$x^ke^{\lambda x}$ for $0\le k\le d_\lambda$. The operator $A_f(\D)$ annihilates
it. Indeed, $(\D-\lambda)$ differentiates the polynomial in the corresponding
block, and all constant-coefficient factors commute. Repeating it $d_\lambda+1$
times kills that block; the other factors do not raise its degree.

This is not the repository's discrete zero-observation closure in new notation.
The finite differential space is used here to propagate an \emph{order
inequality} through positive integrating factors on a continuous domain. A
finite list of derivative identities alone would not establish that inequality.

The representation permits holes in a polynomial's coefficients. The dimension
bound nevertheless uses its largest degree, not the number of nonzero terms:
differentiation can fill those holes. A sparse implementation may store only
nonzero entries, but must not confuse sparse support size with the length of an
annihilator.

\subsection{Source-language boundaries}

A first Lean reifier should recognise real-valued expressions generated by
rational constants, one distinguished real local variable, addition,
subtraction, multiplication, natural powers, and
$\operatorname{Real.exp}(r x)$ with rational $r$. It should produce both the
canonical object and a proof that evaluation agrees with the source expression.
An expression hash is useful for caching but is not that proof.

Several natural-looking extensions require additional work. A shifted argument
$e^{r x+b}$ factors as $e^b e^{rx}$, but $e^b$ need not be rational; it is outside
\eqref{eq:fragment} unless a positive common factor can be removed or the
coefficient domain is deliberately enlarged. Likewise $2^x$ has rate $\log2$,
which is not a rational rate. The analytic ladder theorem allows arbitrary real
factors, but the delivered search and canonical arithmetic are rational. The
prototype does not pretend that a rational approximation to $\log2$ gives an
exact source translation.

Expressions such as $e^{x^2}$, $\sin x$, nested exponentials, and general
multivariate exponential sums are not accepted by this reifier. Some can be
reduced to the fragment by a proved substitution, but there is no implicit
fallback which treats an unsupported operation as an unrelated formal symbol.

\subsection{Why the origin is an especially good anchor}

At $x=0$, every exponential is one and all positive powers of $x$ vanish:
\[
 f(0)=\sum_{\lambda}P_\lambda(0)\in\Q.
\]
Every member of a derivative ladder still has the same form. Thus every initial
sign at the origin is an exact rational comparison, including equality to zero.
No precision loop is needed to certify a multiple boundary zero.

At a rational anchor $a\ne0$, the coefficients and operators remain rational,
but initial values are finite sums of rational multiples of $e^{\lambda a}$.
They require proved enclosures, or a separate symbolic cancellation argument.
The implementation evaluates $P_\lambda(a)$ exactly before approximating any
exponential. Canonical cancellation at this stage is essential: a genuinely
zero value must not be rediscovered by trying to squeeze an interval of
uncertain sign to an impossible finite width.
